\section{Antecedentes}
Cloud Computing es una plataforma esencial y necesaria para gestionar, analizar y guardar el Big Data favor a su potencia elevada de escalar y su flexibilidad de procesar una gran cantidad de información no estructurada donde representa la columna vertebral que soporte esa gran cantidad de información. La integración de Big data con la nube permite a las organizaciones, los Bancos, las empresas y muchas más entidades reducir muchos costos operativos por la migración infraestructura física a un modelo digital adaptado a sus necesidades que se basa en el principio de pago por uso(pay-as-yo-go) [Big data analytics in Cloud computing: an overview]. La inteligencia artificial transforma la nube desde una herramienta tradicional de almacenamiento a un sistema inteligente que permite automatizar los procesos complicados y mejorar la gestión de los recurses automáticamente, transformando los entornos de la nube en sistemas activos capaces de optimizar el rendimiento de forma autónoma. [Applications of AI in Cloud Computing]. Las tecnologías de inteligencia artificial integradas con la nube contribuyen a proporcionar perspectivas predictivas precisas que ayudan las empresas anticipar los desafíos y mejorar sus decisiones según los patrones generados del Big data lo que cual aumenta el valor operativo de los datos en bruto [The Future of AI in Big Data]. El Big data en la nube se enfrenta a una serie de desafíos críticos y de seguridad relacionados con la privacidad y la protección de datos lo que hace necesario es encontrar mecanismos eficientes, avanzados y segura de análisis que garantiza la integridad de la información sensible [Artificial Intelligence Data Security Evaluation]. Por ello, utilizando algoritmos de AI para analizar y evaluar la protección de datos y descubrir las vulnerabilidades y las amenazas cibernéticas de manera proactiva en los varios entornos de la Big data en la nube lo que ofrece una capa inteligente que predice los riesgos antes que ocurran [Artificial Intelligence Data Security Evaluation]. El futuro de la nube inteligente dirige hacia una integración más profunda de inteligencia artificial en la infraestructura de la nube lo que refuerza su capacidad de la gestión de las tareas más complejas y sostenibles [The Future of AI in Big Data]. No obstante, este gran avance tecnológico, surge una clara burbuja de investigación de como equilibrar estos sistemas inteligentes entre su rendimiento alto y la sostenibilidad energética y con el aumento exponencial de la cantidad de datos y el consumo energético de los centros de datos se aparece la necesidad imperiosa de innovar modelos que mejoran la optimización de rendimiento tradicional y se enfocan especialmente en la predicción proactiva de los patrones de carga( workload pattrens) para potenciar la eficiencia energética. Y eso es el núcleo fundamental que busca mi estudio, atreves de subsanar la deficiencia en los estudios actuales que se enfocan en el rendimiento y la seguridad y se ignoran los efectos ambientales y económicos del consumo energético en los entornos del Big Data en la nube. 
